\par\begingroup
\setlength{\lineskip}{4pt}\setlength{\lineskiplimit}{4pt}
\Needspace{8\baselineskip}\bigskip
\pdfbookmark[3]{Darstellung, 18}{pallas-II-article18}
\begin{center}18.\end{center}
\noindent\hspace*{\parindent}Es sei \(\mathrm{f}t\) eine Funktion von \(t\), deren Zahlenwerte für \(t=T,\ t=T+\theta,\ t=T+2\theta,\ t=T+3\theta\) usw. vorliegen. Wir bezeichnen mit \(\mathrm{f}'t\) die Differenz
\[
\mathrm{f}(t+\tfrac12\theta)-\mathrm{f}(t-\tfrac12\theta),
\]
und entsprechend
\[
\begin{aligned}
\text{mit }\mathrm{f}''t&\text{ die Differenz }\mathrm{f}'(t+\tfrac12\theta)-\mathrm{f}'(t-\tfrac12\theta)\\
\mathrm{f}'''t&\text{ die Differenz }\mathrm{f}''(t+\tfrac12\theta)-\mathrm{f}''(t-\tfrac12\theta)\\
\mathrm{f}^{\mathrm{IV}}t&\text{ die Differenz }\mathrm{f}'''(t+\tfrac12\theta)-\mathrm{f}'''(t-\tfrac12\theta)
\end{aligned}
\]
und so weiter.
\par Durch Bildung der Differenzenreihen erhalten wir die Werte von \SourceLeaf{468}\(\mathrm{f}'(T+\tfrac12\theta),\ \mathrm{f}'(T+\tfrac32\theta),\ \mathrm{f}'(T+\tfrac52\theta)\) usw., \(\mathrm{f}''(T+\theta),\ \mathrm{f}''(T+2\theta),\ \mathrm{f}''(T+3\theta)\) usw., \(\mathrm{f}'''(T+\tfrac32\theta),\ \mathrm{f}'''(T+\tfrac52\theta)\) usw. usw. Nehmen wir ferner an, dass
\[
\begin{gathered}
\begin{aligned}
\mathrm{f}t&=\mathrm{f}^{(-1)}(t+\tfrac12\theta)-\mathrm{f}^{(-1)}(t-\tfrac12\theta)\\
\mathrm{f}^{(-1)}t&=\mathrm{f}^{(-2)}(t+\tfrac12\theta)-\mathrm{f}^{(-2)}(t-\tfrac12\theta)
\end{aligned}\\
\text{usw.}
\end{gathered}
\]
Man kann \(\mathrm{f}^{(-1)}(T-\tfrac12\theta)\) einen beliebigen Wert zuweisen, und durch fortgesetzte Addition der Werte von \(\mathrm{f}t\) erhält man die Werte von \(\mathrm{f}^{(-1)}(T+\tfrac12\theta),\ \mathrm{f}^{(-1)}(T+\tfrac32\theta),\ \mathrm{f}^{(-1)}(T+\tfrac52\theta)\) usw. Ebenso weist man \(\mathrm{f}^{(-2)}(T-\theta)\) einen beliebigen Wert zu, und durch fortgesetzte Addition der Werte von \(\mathrm{f}^{(-1)}t\) erhält man die Werte von \(\mathrm{f}^{(-2)}T,\ \mathrm{f}^{(-2)}(T+\theta),\ \mathrm{f}^{(-2)}(T+2\theta)\) usw. und so weiter.
\par Unter diesen Voraussetzungen erhält man
\[
\begin{aligned}
\mathrm{I.}\quad&\int\mathrm{f}t.\,\mathrm{d}t=\alpha+\theta\bigl\{\mathrm{f}^{(-1)}t+\tfrac1{24}\mathrm{f}'t\\
&\qquad-\tfrac{17}{5760}\mathrm{f}'''t+\tfrac{367}{967680}\mathrm{f}^{\mathrm{V}}t\\
&\qquad-\tfrac{27859}{464486400}\mathrm{f}^{\mathrm{VII}}t+\tfrac{1295803}{122624409600}\mathrm{f}^{\mathrm{IX}}t-\text{usw.}\bigr\}
\end{aligned}
\]
\[
\begin{aligned}
\mathrm{II.}\quad&\int\!\int\mathrm{f}t.\,\mathrm{d}t^2=\alpha t+\beta+\theta\theta\bigl\{\mathrm{f}^{(-2)}t+\tfrac1{12}\mathrm{f}t\\
&\qquad-\tfrac1{240}\mathrm{f}''t+\tfrac{31}{60480}\mathrm{f}^{\mathrm{IV}}t\\
&\qquad-\tfrac{289}{3628800}\mathrm{f}^{\mathrm{VI}}t+\tfrac{317}{22809600}\mathrm{f}^{\mathrm{VIII}}t-\text{usw.}\bigr\},
\end{aligned}
\]
wobei \(\alpha\) und \(\beta\) konstante Größen sind.
\par Die Koeffizienten der ersten Reihe erhält man durch Division von Eins durch
\[
\begin{aligned}
&1-\tfrac1{24}xx+\tfrac3{640}x^4-\tfrac5{7168}x^6+\text{usw.}\\
&\quad=1-\tfrac18\cdot\tfrac13xx+\dfrac{1.3}{8.16}\cdot\tfrac15\cdot x^4\\
&\qquad\quad-\dfrac{1.3.5}{8.16.24}\cdot\tfrac17\cdot x^6+\dfrac{1.3.5.7}{8.16.24.32}\cdot\tfrac19\cdot x^8-\text{usw.}\\
&\quad=\dfrac{\log(\BareRoot{(1+\tfrac14xx)}+\tfrac12x)}{\tfrac12x}\mathpunct{\cdot}
\end{aligned}
\]
Quadriert man den Quotienten, so erhält man die Koeffizienten der zweiten Reihe. Höhere Potenzen des Quotienten würden ebenso die Koeffizienten analoger Reihen liefern, mit denen sich die Integrale \(\int^3\mathrm{f}t.\,\mathrm{d}t^3\) usw. ausdrücken lassen, die wir für unsere Untersuchung nicht benötigen.
\par Hier sind zwei weitere Reihen, die man verwenden kann:
\[
\begin{aligned}
\mathrm{III.}\quad&\int\mathrm{f}t.\,\mathrm{d}t=\alpha+\tfrac12\theta\bigl\{\\
&\quad(\mathrm{f}^{(-1)}(t-\tfrac12\theta)+\mathrm{f}^{(-1)}(t+\tfrac12\theta))\\
&\quad-\tfrac1{12}(\mathrm{f}'(t-\tfrac12\theta)+\mathrm{f}'(t+\tfrac12\theta))\\
&\quad+\tfrac{11}{720}(\mathrm{f}'''(t-\tfrac12\theta)+\mathrm{f}'''(t+\tfrac12\theta))\\
&\quad-\tfrac{191}{60480}(\mathrm{f}^{\mathrm{V}}(t-\tfrac12\theta)+\mathrm{f}^{\mathrm{V}}(t+\tfrac12\theta))+\text{usw.}\bigr\}
\end{aligned}
\]
\[
\begin{aligned}
\mathrm{IV.}\quad&\int\!\int\mathrm{f}t.\,\mathrm{d}t^2=\alpha t+\beta+\tfrac12\theta\theta\bigl\{\\
&\quad(\mathrm{f}^{(-2)}(t-\tfrac12\theta)+\mathrm{f}^{(-2)}(t+\tfrac12\theta))\\
&\quad-\tfrac1{24}(\mathrm{f}(t-\tfrac12\theta)+\mathrm{f}(t+\tfrac12\theta))\\
&\quad+\tfrac{17}{1920}(\mathrm{f}''(t-\tfrac12\theta)+\mathrm{f}''(t+\tfrac12\theta))\\
&\quad-\tfrac{367}{193536}(\mathrm{f}^{\mathrm{IV}}(t-\tfrac12\theta)+\mathrm{f}^{\mathrm{IV}}(t+\tfrac12\theta))+\text{usw.}\bigr\}.
\end{aligned}
\]
\SourceLeaf{469}Die Koeffizienten der dritten Reihe erhält man, indem man \(1-\tfrac18xx+\dfrac{1.3}{8.16}x^4-\dfrac{1.2.5}{8.16.24}x^6+\text{usw.}\) durch \(1-\tfrac18\cdot\tfrac13xx+\dfrac{1.3}{8.16}\cdot\tfrac15x^4-\dfrac{1.3.5}{8.16.24}\cdot\tfrac17x^6+\text{usw.}\) dividiert oder den folgenden Ausdruck in eine Reihe entwickelt
\[
\dfrac{\tfrac12x}{\BareRoot{(1+\tfrac14xx)}.\log(\BareRoot{(1+\tfrac14xx)}+\tfrac12x)}\mathpunct{\cdot}
\]
Ebenso ergeben sich die Koeffizienten der vierten Reihe aus der Entwicklung von
\[
\dfrac{\tfrac14xx}{\BareRoot{(1+\tfrac14xx)}.\{\log(\BareRoot{(1+\tfrac14xx)}+\tfrac12x)\}^2}\mathpunct{\cdot}
\]
Noch einfacher erhält man die Koeffizienten der vierten Reihe, indem man die der ersten mit \(1,\ -1,\ -3,\ -5,\ -7\) usw. multipliziert. Auch Integrale höherer Ordnung ließen sich auf ähnliche Weise ausdrücken.
\par Wir lassen hier die Beweise dieser Sätze weg, da sie sich leicht aus der Theorie der erzeugenden Funktionen ableiten lassen, die ein namhafter Geometer gerade dargelegt hat.
\par\endgroup
